\section{Introduction}
\label{sec:introduction}

Large language models are now widely used in bilateral settings such as chat assistants and individual copilots \citep{Maedche2019,Nah2023}. Yet many real work settings require digital support for several humans who coordinate through ongoing communication and maintain a shared operational picture rather than seek isolated answers. Recent work on human-AI collaboration therefore argues for moving beyond assistant-centric interaction toward AI systems that support team processes and shared coordination \citep{Seeber2020,Anthony2023,Bankins2024}. In mission-critical domains, one concrete support need is to maintain an explicit view of what has been assigned, what remains pending, and what has been completed.

Firefighter incident command provides a concrete instance of this broader problem. Commanders coordinate multiple units through short, interleaved radio transmissions while tracking operational tasks. We study this setting through a command-support dashboard whose task list is already derived from procedures and incident context. The model does not discover new tasks or make operational decisions. Instead, it proposes conservative state updates for predefined tasks under human oversight.

Once such a dashboard is the target system, a more specific technical question becomes central: how much transcript structure is actually needed for LLM-based task monitoring from multi-party radio traffic? Rich transcript structure may be costly or brittle to obtain in this setting. Speaker identity may require diarization or prior speaker models, both of which remain difficult under noisy multi-speaker conditions and heterogeneous channels \citep{Mehrish2023A}. Utterance boundaries may be supplied by voice activity detection, but pause-based segmentation can split incomplete thoughts, while more refined speaker-change or semantic segmentation adds models, engineering effort, and delay \citep{Mehrish2023A}. For a command-support system, these are not cosmetic transcript properties but potential deployment bottlenecks because real-time speech systems must trade off model complexity against latency and computational load \citep{Michelsanti2020An}.

This motivates the central research question of the paper: \emph{how much transcript structure is actually needed to track predefined task state reliably from ongoing multi-party communication in this operational setting?} More concretely, we ask whether task-state monitoring depends on speaker identities and explicit utterance boundaries, or whether a simpler transcript representation suffices for this monitoring task.

We address this question through a controlled evaluation of closed-world task-state monitoring across transcript conditions that vary speaker identity and utterance boundaries while keeping the underlying scenario content fixed. The paper contributes in three ways. First, we formulate a bounded coordination-support problem centered on monitoring the state of known operational tasks in firefighter incident command. Second, we introduce a source-grounded synthetic dataset construction and validation pipeline that enables controlled comparison of transcript-structure effects while holding operational content constant. Third, we report results for incremental monitoring and a secondary offline full-transcript reference in a reproducible evaluation pipeline with auditable artifacts and deterministic tests\footnote{GitHub repository with complete data, code and results not revealed due to anonymisation}.

% Update footnote to: https://github.com/zhaw-iwi/swisstext26_pub

The remainder of the paper is structured as follows. Section~\ref{sec:background} situates the work in the literature on multi-human AI support and facilitation. Sects.~\ref{sec:application}--\ref{sec:experimental-setup} then cover the application framing, dataset, and experimental design, before Sect.~\ref{sec:evaluation} reports the results and Sect.~\ref{sec:conclusion} concludes.
